% ===========================================================================
\subsection{\file{.github/workflows-disabled/deploy.yml.disabled}}
% ===========================================================================

\textbf{YAML, 103 lines.} The automation that used to publish this site: it watched
for pushes, checked for a marker in the commit message, copied the files to the
agency's server over an encrypted connection, then checked that the site still
answered. It is currently switched off.

\begin{provenance}{.github/workflows-disabled/deploy.yml.disabled}
Hand-written, added in commit \code{05210af} on 2026-07-28 at the path
\file{.github/workflows/deploy.yml}, rewritten in \code{692e162} the same day when
the deploy target changed from Netlify to the VPS, fixed four times on 2026-07-29,
and finally renamed to its present path in \code{0dd839e}.

It is \textbf{nearly identical to the sibling repository's workflow}: comparing the
two files directly shows differences on exactly six lines, all of them the site's
own name or path (the header comment, the serving path, the workflow name, the
\code{rsync} destination and the verification URL). Every mechanism, every retry
loop and every flag is the same.

This repository is where the evidence for those mechanisms lives. Its history
contains two commits explaining the retry loops by name (\code{6b1aa25},
\code{6fcffee}); the sibling has the loops and no such commits. The
\code{[deploy]} convention itself is workspace-wide and predates both.
\textbf{Confidence: certain for hand-written and for the rename history, both
recorded in git; certain for the six-line difference, checked directly.}
\end{provenance}

\subsubsection*{Why it exists}

To remove a manual step. Without it, publishing means somebody connecting to the
server and copying files by hand, which is slow and easy to get wrong. With it,
publishing is a consequence of writing a commit message in a particular way.

\begin{jargon}{YAML}
A file format for configuration where structure is expressed by indentation rather
than brackets. Two spaces of indent mean ``this belongs to the line above''. It is
unforgiving about indentation, which is its main drawback: one misplaced space
changes the meaning.
\end{jargon}

\begin{jargon}{GitHub Actions, and a workflow}
A service running jobs on GitHub's computers when something happens in a
repository. The jobs are described in a workflow file in \file{.github/workflows/}.
Each job runs on a fresh, empty machine, which is why the file below installs its
own SSH key every time.
\end{jargon}

\subsubsection*{Line by line}

\paragraph{Lines 1 to 19: the header comment and the trigger.}

\begin{lstlisting}[firstnumber=1,caption={\file{deploy.yml.disabled}, lines 1--19}]
# salmanadnan.com: deploy to VPS when commit subject starts with [deploy].
#
# A push to main always checks the marker. It DEPLOYS only when the commit
# subject line STARTS with [deploy], matching every other agency repository.
# Mentioning the marker anywhere else in the message does nothing.
#
# The site is served by Caddy from /root/salmanadnan.com on the VPS.
# Caddy picks up the new files immediately, no restart needed.
#
# Required repository secrets:
#   VPS_HOST     46.202.194.85
#   VPS_USER     deploy
#   VPS_SSH_KEY  private key for that user
name: Deploy salmanadnan.com

on:
  push:
    branches: [main]

\end{lstlisting}

Lines 1 to 13 are a comment block, and it is the best documentation in the
repository. In YAML a \code{\#} starts a comment. It states the marker rule twice
and in capitals, which tells you it had caught somebody out; it names the serving
path and notes that Caddy needs no restart; and it lists the three secrets required.

\begin{honest}{The server's address is written in the comment explaining it is a
secret}
Line 11 records \code{VPS\_HOST 46.202.194.85} in plain text immediately under the
heading ``Required repository secrets''. Whatever the value of keeping the address
in GitHub's secret store, writing it here defeats it.

Proportion matters: an IP address is not a credential, this repository is private,
and the machine is publicly reachable anyway. This document confirmed independently
that salmanadnan.com resolves to exactly this address today, which incidentally
makes the comment currently accurate. The genuine cost is that it is permanent in
the git history, so it becomes misleading rather than merely redundant if the server
is ever replaced. The private key, which is what actually matters, is correctly
absent.
\end{honest}

Line 14 names the workflow, the label shown in GitHub's interface.

Lines 16 to 18 set the trigger: every push to \code{main}, and nothing else. Not on
pull requests, not on a schedule, not manually. That last absence is worth noting:
there is no \code{workflow\_dispatch}, so there is no way to publish without making
a commit.

\paragraph{Lines 20 to 38: the marker check.}

\begin{lstlisting}[firstnumber=20,caption={\file{deploy.yml.disabled}, lines 20--38}]
jobs:
  check-deploy:
    name: Check Deploy Marker
    runs-on: ubuntu-latest
    outputs:
      should_deploy: ${{ steps.check.outputs.should_deploy }}
    steps:
      - id: check
        env:
          COMMIT_MSG: ${{ github.event.head_commit.message }}
        run: |
          SUBJECT="${COMMIT_MSG%%$'\n'*}"
          if [[ "$SUBJECT" == "[deploy]"* ]]; then
            echo "should_deploy=true" >> "$GITHUB_OUTPUT"
            echo "Deploy marker found on the subject line."
          else
            echo "should_deploy=false" >> "$GITHUB_OUTPUT"
            echo "No deploy marker: skipping deploy."
          fi
\end{lstlisting}

The first job exists only to answer a yes-or-no question, and runs on a fresh Ubuntu
machine to do it. Lines 24 and 25 declare that it produces an output called
\code{should\_deploy}, which the second job will read.

Line 28 passes the commit message in as an environment variable rather than pasting
it into the script text. That distinction is a genuine security measure.

\begin{keyidea}{Why the commit message goes through an environment variable}
If line 31 had been written \code{SUBJECT="\$\{\{
github.event.head\_commit.message \}\}"}, GitHub would paste the commit message
directly into the script before the shell ever saw it. Anyone who could write a
commit message could then write shell commands into it and have them run on the
machine holding the deploy key.

Passing it as an environment variable and reading it as \code{\$COMMIT\_MSG} means
the message is data, never code. This is a well-known class of vulnerability in
GitHub Actions, and this file gets it right.
\end{keyidea}

Line 31 extracts the subject line. \code{\$\{COMMIT\_MSG\%\%\$'\textbackslash n'*\}}
is shell syntax meaning ``delete the longest match of a newline followed by
anything, from the end'', leaving everything before the first newline.

Line 32 is the test. \code{"\$SUBJECT" == "[deploy]"*} compares with a trailing
wildcard, which is why the marker must \emph{begin} the subject.

Lines 33 and 36 write the answer to \code{\$GITHUB\_OUTPUT}, a file GitHub reads to
collect a step's outputs. Lines 34 and 37 print a human-readable explanation into
the log, a small kindness to whoever is wondering why nothing deployed.

\paragraph{Lines 39 to 63: starting the deploy, and the SSH setup.}

\begin{lstlisting}[firstnumber=39,caption={\file{deploy.yml.disabled}, lines 39--63}]

  deploy:
    name: Deploy to VPS
    runs-on: ubuntu-latest
    needs: check-deploy
    if: needs.check-deploy.outputs.should_deploy == 'true'
    steps:
      - uses: actions/checkout@v4

      - name: Authorize SSH
        env:
          KEY: ${{ secrets.VPS_SSH_KEY }}
          HOST: ${{ secrets.VPS_HOST }}
        run: |
          mkdir -p ~/.ssh && chmod 700 ~/.ssh
          printf '%s\n' "$KEY" > ~/.ssh/deploy_key
          chmod 600 ~/.ssh/deploy_key
          for i in 1 2 3 4 5; do
            if ssh-keyscan -T 10 -H "$HOST" >> ~/.ssh/known_hosts 2>/tmp/keyscan.err; then
              [ -s ~/.ssh/known_hosts ] && break
            fi
            echo "ssh-keyscan attempt $i failed: $(cat /tmp/keyscan.err 2>/dev/null)"
            sleep 5
          done
          [ -s ~/.ssh/known_hosts ] || { echo "::error::could not fetch the host key"; exit 1; }
\end{lstlisting}

Line 43 makes this job wait for the first, and line 44 runs it only if the answer was
yes. Everything below is skipped entirely on an ordinary commit.

Line 46 checks out the repository onto the machine using GitHub's own action.
Without this the machine would be empty.

Lines 53 to 55 install the private key: create the \file{.ssh} folder, write the key
into it, and set permissions so only the owner can read it. \code{chmod 600} is not
optional; SSH refuses to use a key file other users can read, and its error message
when it does is famously unhelpful.

\begin{jargon}{SSH key pair, and the host key}
Connecting securely involves two separate keys and confusing them is common. The
\emph{deploy key} here is the private half of a pair proving \emph{we} are allowed
in; its public half sits on the server. The \emph{host key} is the server's own
identity, which the client checks to be sure it is talking to the right machine and
not an impostor. Lines 53 to 55 install the first; lines 56 to 63 fetch the second.
\end{jargon}

Lines 56 to 62 are the retry loop, and this repository records exactly why it exists.
\code{ssh-keyscan} asks the server for its host key with a 10-second timeout,
appending the result to \file{known\_hosts} and any complaint to a temporary file.
\code{[ -s ~/.ssh/known\_hosts ]} then tests that the file is non-empty, because
\code{ssh-keyscan} can succeed as a command while producing nothing at all. That
distinction is the subtle part: checking the exit code alone would not have been
enough.

\begin{keyidea}{Two commits name this problem outright}
Commit \code{6b1aa25} is titled ``Retry after transient keyscan failure'' and
\code{6fcffee} ``Retry, runner IP lottery''. Read together they are a complete
diagnosis: GitHub's runners come from a wide pool of addresses, some of them were
being refused or timing out, and the same push succeeded on a second attempt.

This is the clearest example in the whole pair of repositories of a challenge with a
scar. The sibling has the identical retry loop and no commit explaining it, so
anybody reading only that repository would see defensive code with no stated cause.
\end{keyidea}

Line 63 fails the whole job with a message if five attempts produced nothing.
\code{::error::} is GitHub's syntax for marking a log line as an error so it surfaces
in the interface.

\begin{keyidea}{Why fetching the host key each time is a real trade-off}
Fetching a server's identity at the moment you connect, and trusting whatever comes
back, means you are not verifying anything: an attacker positioned between the
runner and the server could supply their own. The secure alternative is to store the
expected host key as a secret and compare.

That is a real weakness and it is a widespread convention, and another project in
this workspace, \code{passbolt-digitalise}, does it the stricter way by keeping a
\file{.github/known\_hosts} file in the repository. Two approaches to the same
problem inside one workspace is worth knowing about.
\end{keyidea}

\paragraph{Lines 64 to 83: copying the files.}

\begin{lstlisting}[firstnumber=64,caption={\file{deploy.yml.disabled}, lines 64--83}]

      - name: Sync files
        env:
          HOST: ${{ secrets.VPS_HOST }}
          USER: ${{ secrets.VPS_USER }}
        run: |
          retry() {
            for attempt in 1 2 3; do
              "$@" && return 0
              echo "attempt $attempt failed: $*"
              [ "$attempt" -lt 3 ] && sleep $(( attempt * 15 ))
            done
            return 1
          }
          retry rsync -az --no-owner --no-group --omit-dir-times --delete \
            -e "ssh -i ~/.ssh/deploy_key -o BatchMode=yes -o ConnectTimeout=25" \
            --exclude '.git' \
            --exclude '.github' \
            --exclude 'assets/ASSETS.md' \
            ./ "$USER@$HOST:/root/salmanadnan.com/"
\end{lstlisting}

Lines 70 to 77 define a small shell function running whatever it is given up to
three times, waiting longer between attempts: 15 seconds, then 30. \code{"\$@"}
means ``all the arguments, run as a command''. This is the second retry loop, and
like the first it is evidence of a step that failed in practice.

Line 78 is the copy itself. Each flag earns its place:

\begin{table}[H]\centering\small
\begin{tabular}{@{}ll@{}}
\toprule
Flag & What it does \\
\midrule
\code{-a} & archive mode: recurse, and preserve timestamps, permissions and links \\
\code{-z} & compress during transfer \\
\code{--no-owner} & do not try to set file ownership on the far end \\
\code{--no-group} & do not try to set group ownership either \\
\code{--omit-dir-times} & do not try to set directory timestamps \\
\code{--delete} & remove files on the server that no longer exist here \\
\bottomrule
\end{tabular}
\caption{The \code{rsync} flags, and why each is present.}
\end{table}

The three \code{--no} and \code{--omit} flags all work around the same thing: the
\code{deploy} user does not own the target directory and cannot change ownership or
directory timestamps. Without them \code{rsync} reports errors and fails. These
flags are the surviving trace of the problem commit \code{431ebbf} was fixing.

\code{--delete} deserves respect: it makes the server an exact mirror, so removing a
file here removes it there. That is usually what you want, and it is also the flag
that turns a mistaken source path into a wiped site.

Line 79 tells \code{rsync} how to connect: use this key, never prompt for a password
(\code{BatchMode=yes}, essential because there is no human to answer), and give up
after 25 seconds.

Lines 80 to 82 exclude three things: the git history, the workflow folder, and the
internal checklist.

Line 83 is the source and destination. \code{./} means the current directory's
\emph{contents}, and the trailing slash is significant: without it \code{rsync}
would copy the directory itself into the target, producing a nested folder.

\begin{honest}{This line would now take down the live site}
\code{./} was the right source when \file{index.html} sat at the repository root.
Since commit \code{0dd839e} moved everything one level down, \code{./} contains a
single folder called \file{new-design-reference}. Restoring this workflow unchanged
would copy that folder to the server and, because \code{--delete} is set, remove the
real site's files in the same operation.

This is not hypothetical harm. Section 1.2 established that salmanadnan.com is live,
answering with \code{200}, and resolves to the very address on line 11. The correct
source after the move would be \code{new-design-reference/}.

The same commit that moved the files disabled this workflow, so the two are
consistent; the trap exists only for somebody who re-enables one without reading the
other.
\end{honest}

\begin{honest}{The exclusion for the checklist no longer matches}
\code{--exclude 'assets/ASSETS.md'} was correct when that path was correct. The file
now lives at \file{new-design-reference/assets/ASSETS.md}, which this pattern does
not match, so a restored workflow would publish the internal checklist, including
its instruction to drop a photograph in, to the public internet.
\end{honest}

\paragraph{Lines 84 to 103: verifying, and cleaning up.}

\begin{lstlisting}[firstnumber=84,caption={\file{deploy.yml.disabled}, lines 84--103}]

      - name: Verify
        env:
          HOST: ${{ secrets.VPS_HOST }}
          USER: ${{ secrets.VPS_USER }}
        run: |
          for i in $(seq 1 12); do
            code=$(ssh -i ~/.ssh/deploy_key -o BatchMode=yes "$USER@$HOST" \
              'curl -sf -o /dev/null -w "%{http_code}" --max-time 10 https://salmanadnan.com/' || echo 000)
            case "$code" in
              2*|3*) echo "Site answering with $code after ${i} check(s)."; exit 0 ;;
            esac
            echo "waiting, got $code ($i/12)"
            sleep 5
          done
          echo "::warning::site did not respond 2xx/3xx; check Caddy on the VPS"

      - name: Clean up key
        if: always()
        run: rm -f ~/.ssh/deploy_key
\end{lstlisting}

The verification step asks, up to twelve times over about a minute, whether the site
is answering. It does this by connecting to the server and running \code{curl}
\emph{there} rather than from the runner, which tests the server's own view of
itself.

The \code{curl} flags: \code{-s} quiet, \code{-f} treat an error status as a
failure, \code{-o /dev/null} throw the page away,
\code{-w "\%\{http\_code\}"} print just the status number, and a 10-second limit.
The \code{|| echo 000} supplies a fallback so a failed connection yields the harmless
string \code{000} rather than an empty one.

Line 93 accepts any status beginning with 2 or 3, meaning success or redirect, and
exits successfully. That tolerance is well judged and this document can confirm it
matters: \code{http://salmanadnan.com/} currently answers \code{308}, a permanent
redirect to the secure version, so a check accepting only \code{200} could fail on a
perfectly healthy site.

\begin{honest}{A failed verification is only a warning}
Line 99 uses \code{::warning::} rather than \code{::error::}, and the step does not
exit with a failure code. So if the site never came back, the workflow would still
finish green with a warning buried in the log.

Given that the preceding step ran \code{rsync --delete} against a live site, this is
the one moment you would want a loud red failure and an email. It is the clearest
correctness weakness in the workflow, and it is one word to fix.
\end{honest}

Lines 101 to 103 delete the private key from the runner. \code{if: always()} makes
this run even when an earlier step failed, which is the point: a key must not be left
behind because something else went wrong. In practice GitHub destroys the machine
anyway, so this is belt and braces, and it is the right instinct.

\subsubsection*{What it connects to}

\textbf{Reads:} three GitHub secrets (\code{VPS\_SSH\_KEY}, \code{VPS\_HOST},
\code{VPS\_USER}) and the commit message of the push that triggered it.

\textbf{Writes:} the contents of \file{/root/salmanadnan.com/} on the VPS at
\code{46.202.194.85}, which Caddy serves at \code{https://salmanadnan.com/}.

\textbf{Excludes:} \file{.git}, \file{.github} and \file{assets/ASSETS.md}.

\textbf{Runs when:} never, currently. It sits in \file{.github/workflows-disabled/}
with the extension \file{.disabled}, and GitHub only reads
\file{.github/workflows/*.yml}.

\subsubsection*{Concepts introduced}

YAML; GitHub Actions jobs, steps, outputs and secrets; environment variables as a
defence against injection; SSH key pairs and host keys; \code{rsync} and its
archive and delete semantics; retry loops with growing backoff; and the
\code{[deploy]} marker convention shared across this workspace.

\subsubsection*{How it breaks}

\begin{itemize}
\item \textbf{Restored unchanged}, it copies the wrong directory and, with
  \code{--delete}, takes down the live salmanadnan.com.
\item \textbf{A missing or wrong \code{VPS\_SSH\_KEY}} fails at the \code{rsync}
  step after three attempts. Commit \code{cfa65d7} fixed this.
\item \textbf{A \code{deploy} user without permission to traverse into
  \code{/root}} fails the same way. Commit \code{431ebbf} fixed this, and it is the
  least obvious of the failures: the user can hold the key, connect successfully, and
  still be unable to reach the target directory.
\item \textbf{GitHub's runner being refused by the server} fails the host-key fetch.
  Commits \code{6b1aa25} and \code{6fcffee} are both this.
\item \textbf{A commit whose subject does not start with \code{[deploy]}} does
  nothing at all, by design, and the log says so.
\item \textbf{The site failing to come back} produces a warning and a green tick,
  which is the failure mode worth fixing first.
\end{itemize}

\subsubsection*{Honest findings for this file}

The server address is recorded in a comment describing it as a secret; the
\code{rsync} source path and the \file{ASSETS.md} exclusion are both stale after the
move, and here the consequence is a live site going down rather than a parked one;
the host key is trusted on first sight rather than pinned, unlike
\code{passbolt-digitalise} in the same workspace; and a failed post-deploy check
produces a warning rather than a failure. Against those: the commit message is
handled safely as data rather than as code, both retry loops are well built and
documented by the commits that caused them, the status check correctly tolerates
redirects, and the key is cleaned up unconditionally.
